\question[5]
Die Klasse Temperatur hat das Attribut celsius. Dem Konstruktor wird der Wert von celsius übergeben.
Implementiere die Klasse Temperatur unter Wahrung des Geheimnisprinzips mit getter- und setter-Methode.
Der Setter lehnt Werte unter -273.15 ab, das Attribut behält dann seinen Wert. Auch der Konstruktor benutzt
den Setter, der Startwert des Attributs ist 0. Folgender Dialog soll möglich sein:
\begin{lstlisting}
t = Temperatur(20)
t.set_celsius(-300)
print(t.get_celsius())
t.set_celsius(-10)
print(t.get_celsius())
u = Temperatur(-500)
print(u.get_celsius())

20
-10
0
\end{lstlisting}
\begin{solutionbox}{5.5cm}
\begin{lstlisting}
class Temperatur:
    def __init__(self, celsius):
        self.__celsius = 0
        self.set_celsius(celsius)

    def get_celsius(self):
        return self.__celsius

    def set_celsius(self, celsius):
        if celsius >= -273.15:
            self.__celsius = celsius
\end{lstlisting}
\end{solutionbox}
